\section*{历史注记}
\phantomsection
\addcontentsline{toc}{section}{历史注记}

（方括号中的数字指本注记末尾的参考文献。）

我们在此不打算重复复数理论与四元数理论发展的全部历史，因为这些理论本质上属于代数；但我们要谈谈复数的几何表示，它在许多方面都是数学史上决定性的一步。

无疑，第一个对复数与平面点之间的一一对应具有清晰概念的是 C. F. 高斯（*），他还把这一思想应用于复数理论，并预见到 19 世纪的分析学家们将对它作出的种种运用。在 17 与 18 世纪，数学家们逐渐确信：使他们能够解所有二次方程的虚数，也可以用来解任意次数的代数方程。18 世纪期间发表了证明这一定理的许多尝试；但撇开那些仅以循环论证为基础的尝试不谈，没有一种尝试能免于严重的异议。高斯在对这些尝试作了详尽的考察、对其缺陷进行了缜密周详的批评之后，在他的就职论文（写于 1797 年，发表于 1799 年）中提出要最终给出一个严格的证明。他采纳了达朗贝尔 {[}在其 1746 年发表的证明中（**）{]}

高斯指出：平面的点 \((a, b)\) —— 使 \(a + ib\) 为多项式 \(\mathrm{P}(x + iy) = \mathrm{X}(x, y) + i\mathrm{Y}(x, y)\) 之根的那些点 —— 就是曲线 X = o 与 Y = o 的交点；通过对这些曲线的定性研究，他接着证明其中一条曲线的一段连续弧连接着由另一条曲线所围成的两个不同区域中的点，并由此断定两曲线相交（{[}1{]}，第~3 卷，第~3 页；另见{[}1 a{]}）。就其清晰性与独创性而言，这一证明比先前的尝试前进了一大步，而且无疑是把纯拓扑推理应用于代数问题的最早例子之一（*）。

在他的论文中，高斯并没有明确地定义平面点与复数之间的对应；至于复数本身，以及两百年来由它们引起的“存在性”问题，他持保留态度，并有意识地把他所有的论证都以只涉及实量的形式表述出来。但是，他的证明中思想的进程，若不预先假定平面点与复数的一种完全自觉的等同，就会是完全不可理解的；而他在同一时期关于数论与椭圆函数（它们也涉及复数）的研究，又强化了这一揣测。他从熟悉虚数的几何概念到亲手由此获得成果的过程，在他（最近才发表的）把复数应用于初等几何问题求解的笔记（{[}1{]}，第~4 卷，第~396 页及第~8 卷，第~307 页）中清楚地显示出来。更为明确的是他 1811 年致贝塞尔的信（{[}1{]}，第~8 卷，第~90-91 页），其中他勾画了复变量函数积分理论的要点：“正如全部实量的整体可以用一条无穷直线来表示，实量与虚量的全体所成的完整区域也可以用一张无穷平面来表示，其中每个点由其横坐标 a 与纵坐标 b 确定，同时表示量 \(a + ib\) 。于是，从 x 的一个值到另一个值的连续过渡沿一条曲线发生，因而可以以无穷多种方式实现\ldots”。

但直到 1831 年，高斯才（就“高斯整数” \(a + ib\) 的引入，其中 \(a\) 与 \(b\) 为整数）如此明确地公开阐述他在这方面的思想（{[}1{]}，第~2 卷，Theoria Residuorum Biquadraticorum，Commentatis Secunda，art. 38，第~109 页，及 Anzeige，第~174 页及以下）。在这之前的若干年里，用几何方式表示复数的想法已被两位默默无闻的业余爱好者各自独立地重新发现，他们两人都或多或少靠自学成才，对数学没有其他贡献，而且两人与他们那个时代的科学界都几乎没有接触。因此，他们的工作面临完全湮没无闻的危险；这正是时间上居先的那一位的遭遇——丹麦人 C. 韦塞尔所写的一本小册子。它发表于 1798 年，构思与表述都很清晰，但直到一个世纪之后才被人从遗忘中拯救出来。第二位，瑞士人 J. 阿尔冈，也可能遭遇同样的命运，他只是碰巧使自己发表于 1806 年的著作在七年后被重新发掘（*）。这项工作在《热尔冈年刊》上引起了一场热烈的讨论，并且在 1820 至 1830 年间，在法国和英格兰成为几位无名作者若干出版物的主题。但是，需要一位大人物名字的威望来结束这些争论并把数学家们聚集到新观点上来；直到 19 世纪中叶，继高斯（上述）在德国的著作、哈密顿与凯莱在英格兰关于超复系的工作、以及最后柯西（**）在法国的支持之后，复数的几何表示才终于被普遍接受——而仅仅几年之后，黎曼的天才便进一步扩大了几何在解析函数理论中的作用，并同时一手创立了拓扑学这门科学。

\par\noindent\textbf{* \(\ast\)}\par

借助角在圆上所截取的弧来测量角，与角的概念本身一样古老，而且已为巴比伦人所知：他们的角的度量单位是度，我们至今仍在使用。巴比伦人只使用介于 o 与 360° 之间的角的测度；这对他们的目的已经足够，因为这些目的主要是把天体在其视轨道上的位置固定在一些确定的点处，并为科学或占星的目的编制这些轨道的表。

在古典时代的希腊几何学家那里，角的概念（欧几里得《几何原本》，第 I 卷，定义 8 与 9）受到的限制甚至更大，因为它只适用于小于两直角的角；而另一方面，由于他们的比例与度量理论建立在对待测量的量任意大的倍数进行比较的基础上，角对他们来说就不能成为可测量的量，尽管他们当然有相等角、一角大于或小于另一角、以及两角之和（只要和不超过两直角）这些概念。正如分数的加法一样，角的测量在他们眼中必定只是一种经验程序，没有科学价值。这种态度在阿基米德关于螺线的令人赞叹的论文（{[}3{]}，第~2 卷，第~1-121 页）中得到很好的说明：由于不能用向径与角成比例来定义螺线，他给出了一种运动学的定义（定义 1，第~44 页；参看命题 12 的叙述，第~46 页），并且如该著作其余部分所示，他成功地从中提取出了角的测度的一般概念所能给予他的一切，倘若他掌握了这一概念的话。至于希腊天文学家，在这一点上如同在许多其他方面一样，他们似乎满足于追随其巴比伦前辈。

这里也如同实数概念的演变（参看第四章的历史注记）一样，希腊科学衰落时期严格精神的松弛带来了向“朴素”观点的回归，而这种观点在某些方面比刻板的欧几里得观点更接近我们自己的观点。于是一位轻率的增补者在欧几里得（《几何原本》第 VI 卷，命题 33）中插入了下面这个著名的命题：“角与它在圆上截取的弧成比例”（*），而一位对该命题“证明”作评注的无名学者毫不迟疑地——当然是毫无根据的——引入了等于一个圆周任意大倍数的弧以及与这些弧相应的角（**）。但即使 16 世纪的韦达，尽管他在发现方程 \(\sin nx = \sin \alpha\) 有多个根时似乎已接近我们现代的角概念，也只得到了对应于小于两直角的角的那些根（{[}4{]}，第~305 页）。直到 17 世纪，这一观点才被最终超越；在牛顿发现 \(\sin x\) 与 \(\cos x\) 的幂级数展开、从而给出这些函数对变量一切值都有效的表达式之后，欧拉联系“虚”数的对数最终表述了角的测度概念的精确观念（{[}5{]}，(1)，第~17 卷，第~220 页）。

当然，用圆弧长度定义角的测度这一经典定义不仅直观，而且本质上是正确的；然而要使它严格，需要曲线长度的概念，即积分学。从所涉及的结构的角度看，这是一个非常迂回的程序，而且，正如我们在正文中所见，可以只借助拓扑群理论的手段达到同一目的；以这种方式，实指数函数与复指数函数显得出自同一来源，即刻画“单参数群”的定理（第五章，§ 3，定理 1）。

